% Fuente íntegra: originales/cuaderno/REORIENTACION_GENEALOGICA_Y_RESULTADOS.md,
% §§2--7. Gestión y procedencia conservadas aparte.
\Needspace{16\baselineskip}
\section{Régularité généalogique, restitution opératorielle et géométrie de l'action}
\label{md:restop:seccion}

\subsection{Production, compatibilité et restitution}
\label{md:restop:jerarquia}

Le phénomène primaire est une dynamique qui conserve la provenance en produisant de nouvelles observations. Trois opérations aux fonctions mathématiques différenciées sont réunies :
\begin{enumerate}
 \item \emph{Production.} La mise à jour APP--TRIT--TPK prolonge l'état et ses régions corrélées. La phase peut revenir tandis que la mémoire progresse.
 \item \emph{Compatibilité.} Les lectures régionales, angulaires, d'action et d'incidence obéissent à des relations communes. Leur variation doit respecter ces relations pour correspondre à un même état.
 \item \emph{Restitution.} La conservation de l'information complémentaire permet de retrouver les valeurs antérieures, les produits d'opérateurs et les différences entre ordres de composition.
\end{enumerate}
La comparaison entre procédures pour obtenir $\pi$ doit considérer ces trois niveaux. La thèse spécifique de HMT réside dans la construction conjointe et ses relations, outre la production de développements numériques. Une coïncidence décimale ou la périodicité d'une phase ne décrivent pas à elles seules la portée de la construction. Cette distinction identifie l'objet mathématique pertinent pour une comparaison historique ; le développement focal n'établit pas à lui seul une priorité mondiale.

\subsection{Régimes de vacance et retour de la phase de capacité}
\label{md:restop:vacancias}

\subsubsection{Calendrier de capacité et sélecteur tritique}
\label{md:restop:antecedentes-capacidad}

Le lecteur de capacité reçoit la compactification interne $729/1000$ et
distingue l’indice de raffinement de l’indice de capacité. Avec la
convention de la section~\ref{vac-capacidad-trit},
\begin{equation}
 \delta=\log_{10}(10/9),\qquad
 p_j=\left\lfloor\frac{j}{\delta}\right\rfloor,\qquad
 v_j=p_j-j,\qquad j\geq1.
 \label{md:restop:capacidad}
\end{equation}
Les intervalles entre événements sont
\begin{equation}
 h_j=p_{j+1}-p_j\in\{21,22\},\qquad
 s_j=22-h_j.
 \label{md:restop:intervalos}
\end{equation}
La classe tritique obéit à
\begin{equation}
 [v_{j+1}]_3=[v_j-s_j]_3.
 \label{md:restop:clase-tritica}
\end{equation}
Un intervalle court change la classe ; un intervalle long la conserve.
Les séparations $6/7$ décrivent les plateaux de ce régime induit.
La construction maintient son type de facteur d’observation de l’état TPK ;
elle ne remplace pas cet état et n’identifie pas ces séparations à des signes de
charge ou à des chiffres de $\pi$ ou de $e$.
Le sélecteur canonique est
$M_{\mathrm{ph}}(1+[v_j]_9)$, dont le motif dépend de la classe modulo trois.

\subsubsection{Restitution du régime local et déplacement modulo neuf}
\label{md:restop:consecuencia-capacidad}

Soit
\begin{equation}
 \beta_{\mathrm{cap}}=22-\delta^{-1},\qquad
 q_k=\left\lfloor\frac{k}{\beta_{\mathrm{cap}}}\right\rfloor,\qquad k\geq1.
 \label{md:restop:pendiente-inducida}
\end{equation}
Les positions $q_k$ sont les indices des intervalles courts. En effet,
\begin{equation}
 p_j=22j-\lceil j\beta_{\mathrm{cap}}\rceil,\qquad
 s_j=\lceil(j+1)\beta_{\mathrm{cap}}\rceil
     -\lceil j\beta_{\mathrm{cap}}\rceil.
 \label{md:restop:eventos-cortos}
\end{equation}
L’irrationalité de $\delta$ s’obtient à partir des valuations de $2,3,5$
dans l’égalité qu’imposerait $(10/9)^q=10^p$ ; elle implique celle de
$\beta_{\mathrm{cap}}$. Ainsi, les indices positifs ne présentent pas
d’ambiguïté aux extrémités. La discontinuité correspondant à l’entier $k$
se produit en $j=\lfloor k/\beta_{\mathrm{cap}}\rfloor$.

\begin{proposition}
\label{md:restop:proposicion-regimen}
Entre les événements d’indices $q_k$ et $q_{k+3}$, en prenant les intervalles dont
les indices appartiennent à $[q_k,q_{k+3})$, la classe TRIT revient, tandis que la phase
de capacité modulo neuf change.
\end{proposition}
\begin{proof}
Il y a exactement trois intervalles courts dans ce segment.
Les bornes $1/7<\beta_{\mathrm{cap}}<3/20$ donnent
$20<3/\beta_{\mathrm{cap}}<21$, de sorte que
\begin{equation}
 m_k=q_{k+3}-q_k\in\{20,21\}.
 \label{md:restop:duracion-tritica}
\end{equation}
La somme des longueurs et le gain de capacité sont, respectivement,
\begin{equation}
 \Delta p=22m_k-3\in\{437,459\},\qquad
 \Delta v=21m_k-3\in\{417,438\}.
 \label{md:restop:ganancias}
\end{equation}
Dans les deux cas, $\Delta v\equiv0\pmod3$. En revanche,
\begin{equation}
 \boxed{\Delta v\equiv3\ \text{ou}\ 6\pmod9.}
 \label{md:restop:desplazamiento-nonadico}
\end{equation}
L’énoncé est démontré.
Les bornes de $\beta_{\mathrm{cap}}$ disposent de certificats entiers :
\begin{equation}
 \begin{aligned}
 10^{146}>9^{153}&\ \Longrightarrow\
 \delta>\frac7{153},\\
 10^{417}<9^{437}&\ \Longrightarrow\
 \delta<\frac{20}{437}.
 \end{aligned}
 \label{md:restop:certificados-cotas}
\end{equation}
La monotonie de $22-1/x$ produit
$1/7<\beta_{\mathrm{cap}}<3/20$. Ainsi, les bornes reposent sur des
inégalités entières et non sur une substitution décimale arrondie.
\end{proof}

La valeur du sélecteur
$\tau_j=M_{\mathrm{ph}}(1+[v_j]_9)$ revient parce qu’elle dépend de la classe
modulo trois. Ce retour conserve sa distinction avec l’état TRIT
complet et avec un tour chronologique de $\Gamma_9$.
Les deux durées apparaissent une infinité de fois par la densité de la rotation
irrationnelle de pente $1/\beta_{\mathrm{cap}}$.

La régularité présente des niveaux distincts : le régime local est restauré
tandis que la phase de capacité reste décalée. Le retour nonadique de
l’état possède, en outre, son propre accroissement de mémoire. Chaque objet
conserve donc sa notion de retour et sa durée.
Les nombres $437$ et $459$ sont des durées de cette observation induite,
sans qu’on leur attribue le statut de constantes universelles supplémentaires.
Le résultat réunit explicitement les formules précédentes ; sa présentation
comme conséquence focale ne constitue pas une affirmation de priorité
historique par rapport à des développements qui n’ont pas été comparés exhaustivement.

\subsection{Participation de la mémoire à la composition opératorielle}
\label{md:restop:memoria}

\subsubsection{Représentation élargie de lecture et de mémoire}
\label{md:restop:estructura-partida}

La représentation de lecture et de mémoire s’obtient en séparant huit
composantes et une neuvième. Dans une même fibre, avec l’identité comme transport de
référence, le transport interne $C$ induit
\begin{equation}
 T_C=\frac{8I+C}{9},\qquad
 D_C=\frac{\sqrt8(C-I)}{9},\qquad
 R_C=\frac{I+8C}{9},
 \label{md:restop:bloques}
\end{equation}
\begin{equation}
 \mathcal U_C=
 \begin{pmatrix}
 T_C&D_C\\D_C&R_C
 \end{pmatrix}.
 \label{md:restop:transporte-ampliado}
\end{equation}
Le couplage provient du changement de base de la décomposition nonadique.
Pour les transports qui conservent la forme déclarée, il conserve la forme
élargie. L’identité de composition est
\begin{equation}
 \mathcal U_{C_2}\mathcal U_{C_1}
 =\mathcal U_{C_2C_1},\qquad
 T_2T_1+D_2D_1=\frac{8I+C_2C_1}{9},
 \label{md:restop:composicion-ampliada}
\end{equation}
avec $T_i=T_{C_i}$ et $D_i=D_{C_i}$.
Le second terme correspond à la mémoire produite lors de la première
opération et réintroduite dans la seconde. L’écarter change la
procédure : cela produit $T_2T_1$ au lieu de la lecture de la composition
élargie.

\subsubsection{Répartition de la sensibilité à l’ordre de composition}
\label{md:restop:conmutadores}

\begin{proposition}
\label{md:restop:proposicion-orden}
Avec la convention $[A,B]=AB-BA$,
\begin{equation}
 [T_2,T_1]=\frac1{81}[C_2,C_1],\qquad
 [D_2,D_1]=\frac8{81}[C_2,C_1].
 \label{md:restop:reparto-conmutadores}
\end{equation}
Par conséquent, la différence de lecture entre les deux ordres, avec la
mémoire réintroduite, est
\begin{equation}
 \boxed{(T_2T_1+D_2D_1)-(T_1T_2+D_1D_2)
 =\frac19[C_2,C_1].}
 \label{md:restop:orden-completo}
\end{equation}
\end{proposition}
\begin{proof}
En développant les produits de $T$, les termes constants et linéaires
s’annulent dans la différence ; il reste le commutateur divisé par $81$.
En développant les produits de $D$, la même annulation se produit et le
facteur est $8/81$. Leur somme est $1/9$.
\end{proof}

Le mécanisme de réentrée fait partie de la représentation élargie.
La comparaison des deux ordres explicite une conséquence :
dans cette représentation, la contribution de la mémoire à la réponse
antisymétrique est huit fois celle retenue par la composition de lectures
isolées. La réponse complète est neuf fois cette dernière. Les
proportions correspondent à des opérateurs de réponse, non à des pourcentages
d’énergie.

La restitution conserve la non-commutativité et son intensité de lecture.
Dans ce cas, la compression isolée atténue un commutateur non nul sans
l’annuler. Une application à des forces physiques conserve en outre les opérateurs
d’interaction et la fonctionnelle de travail correspondante ; l’identité
algébrique ne les identifie pas à elle seule.

\subsubsection{Confrontation opératorielle et normalisation nonadique}
\label{md:restop:contraste}

Pour le motif algébrique général $(n-1)+1$, les mêmes calculs donnent les
facteurs $1/n^2$, $(n-1)/n^2$ et $1/n$. L’architecture nonadique fixe
$n=9$. On distingue ainsi la loi de composition de la valeur concrète de
sa normalisation.

Lorsque les opérations effectives d’une application sont représentées par
$C_1,C_2$, les procédures qui écartent ou réintroduisent la mémoire
prédisent des réponses différentes à l’ordre. La confrontation porte sur la
composition opérationnelle et sur l’accessibilité de la mémoire, en conservant
les coefficients de la représentation.

\subsubsection{Normes de lecture et de mémoire dans l’inversion centrale}
\label{md:restop:treinta-y-dos}

Pour $C=-I$, on obtient
\begin{equation}
 T^*T=\frac{49}{81}I,\qquad
 D^*D=\frac{32}{81}I.
 \label{md:restop:normas}
\end{equation}
Le numérateur $32$ provient de $8|{-1}-1|^2$.
Dans les neuf composantes de la mémoire, la norme comprend
$8^2+8=72$, avec le dénominateur $27^2$. Le carré $64$ participe à
ce calcul avec les huit autres contributions.
Une identification à une région de $64$ cases exige l’application
d’incidence de la région et ne s’obtient pas en doublant le numérateur $32$.
La lecture géométrique d’un support $8\times8$ par rapport à un autre
$9\times9$ conserve donc son problème de connexion avec
l’application concrète ; l’égalité numérique isolée ne le remplace pas.

\clearpage
\subsection{Caractère d’action et coordonnée complète de réciprocité}
\label{md:restop:h5}

\subsubsection{Coefficients généalogiques et continuation régionale}
\label{md:restop:coeficientes}

La matrice centrale et ses invariants sont
\begin{equation}
 B_c=\begin{pmatrix}7&2\\2&7\end{pmatrix},\qquad
 \lambda_+=9,\qquad\lambda_-=5.
 \label{md:restop:matriz-central}
\end{equation}
Interviennent en outre l’orbite de longueur quatre du pas trois, le
calendrier de longueur douze et le quotient de recensement $104976/1944=54$.
Les relations entre coefficients sont
\begin{equation}
 \begin{aligned}
 \frac9{16}&=\frac{\lambda_+}{4^2},&
 \frac59&=\frac{\lambda_-}{\lambda_+},\\
 \frac7{48}&=\frac{\operatorname{tr}B_c/2}{4\cdot12},&
 \frac1{54}&=\frac{1944}{104976}.
 \end{aligned}
 \label{md:restop:coeficientes-h5}
\end{equation}
L’attribution des degrés et des signes appartient à la définition du caractère ;
leurs cardinaux isolés ne la déterminent pas. Le lecteur complet est
\begin{equation}
 H_5=\frac12\sqrt{\frac{1000\alpha}{\varphi}}-\alpha
 +\frac9{16}\alpha^2-\frac59\alpha^3
 +\frac7{48}\alpha^4-\frac1{54}\alpha^5,
 \label{md:restop:h5-completo}
\end{equation}
\begin{equation}
 \begin{aligned}
 D_A&=\exp(-100\pi\alpha/9),\\
 \mathcal C_\pi&=169\alpha^6+
 \frac{D_A\alpha^7}{1-D_A\alpha},\\
 \eta_{\mathrm{ret}}&=H_5-\frac{90}{\pi}\mathcal C_\pi.
 \end{aligned}
 \label{md:restop:retorno-regional}
\end{equation}
La fraction conserve la continuation infinie définie par
\begin{equation}
 \mathscr R(z)=D_Az^7+D_Az\mathscr R(z),
 \label{md:restop:renovacion}
\end{equation}
avec un poids initial unitaire. La continuation demeure intégrale dans le
lecteur ; un polynôme tronqué ne remplace pas cet objet. La construction
des coefficients et le renouvellement sont réunis dans la
section~\ref{sec:revision-planck}.

\subsubsection{Classification de la réciprocité radiale orientée}
\label{md:restop:geometria}

On conserve $A_{\deg}=1000\alpha$,
$C^*_{\deg}=2(\eta_{\mathrm{ret}}+\alpha)$ et le récupérateur radial du
théorème~\ref{rec:thm:accion}. Pour maintenir les deux orientations, soient
\begin{equation}
 \epsilon\in\{+1,-1\},\qquad
 q_{\mathrm{rad}}=R_\star^4>0,\qquad
 \Xi_{\mathrm{rad}}=
 \epsilon\frac{1-q_{\mathrm{rad}}}{1+q_{\mathrm{rad}}}.
 \label{md:restop:par-radial}
\end{equation}
Le récupérateur \eqref{rec:eq:eta-orientada} s’écrit
\begin{equation}
 \frac{\eta_{\mathrm{ret}}}{\alpha}=500\Xi_{\mathrm{rad}}-1.
 \label{md:restop:recuperador}
\end{equation}
L’involution
$(\epsilon,q_{\mathrm{rad}})
 \mapsto(-\epsilon,q_{\mathrm{rad}}^{-1})$
conserve cette action. La coordonnée qui réunit cette structure est
\begin{equation}
 \boxed{\chi_{\mathrm{rad}}=\epsilon\log q_{\mathrm{rad}}.}
 \label{md:restop:chi}
\end{equation}

\begin{proposition}
\label{md:restop:proposicion-orbitas}
$\chi_{\mathrm{rad}}$ classe exactement les orbites de cette involution
sur $\{\pm1\}\times\mathbb R_{>0}$.
\end{proposition}
\begin{proof}
Changer simultanément le signe et inverser $q_{\mathrm{rad}}$ conserve
$\chi_{\mathrm{rad}}$. Si deux paires ont la même valeur de
$\chi_{\mathrm{rad}}$, avec le même signe, elles ont le même
$q_{\mathrm{rad}}$ ; avec des signes opposés, leurs rayons quartiques sont
réciproques. Ce sont exactement les deux cas d’appartenance à une même
orbite. L’argument comprend $q_{\mathrm{rad}}=1$, dont l’orbite conserve
les deux orientations.
\end{proof}

L’identité
$(1-q)/(1+q)=-\tanh(\log q/2)$ donne
\begin{equation}
 \Xi_{\mathrm{rad}}
 =-\tanh(\chi_{\mathrm{rad}}/2),\qquad
 \boxed{\eta_{\mathrm{ret}}
 =\alpha[-500\tanh(\chi_{\mathrm{rad}}/2)-1].}
 \label{md:restop:accion-chi}
\end{equation}
La compatibilité avec le caractère complet prend la forme
\begin{equation}
 \boxed{\chi_{\mathrm{rad}}
 =-2\operatorname{artanh}
 \left(\frac{H_5+\alpha-(90/\pi)\mathcal C_\pi}{500\alpha}\right).}
 \label{md:restop:chi-caracter}
\end{equation}
Le terme $-\alpha$ de $H_5$ s’annule avec le $+\alpha$ du
contraste angulaire avant cette lecture. Le dénominateur de la
continuation régionale demeure complet. L’égalité reçoit les
sorties et les lecteurs HMT antérieurs et réunit leurs deux représentations ;
elle ne définit pas rétroactivement $\alpha$.

Dans la branche positive d’action et de rayon fini,
\begin{equation}
 0<\eta_{\mathrm{ret}}<499\alpha,\qquad
 \chi_{\mathrm{rad}}
 <-2\operatorname{artanh}(1/500).
 \label{md:restop:dominio-positivo}
\end{equation}
Avec $\alpha$ fija,
\begin{equation}
 \frac{\partial\eta_{\mathrm{ret}}}{\partial\chi_{\mathrm{rad}}}
 =-250\alpha\,\operatorname{sech}^2(\chi_{\mathrm{rad}}/2)<0.
 \label{md:restop:derivada-accion}
\end{equation}
Par conséquent, la section adimensionnelle d’action détermine de manière unique la
classe radiale orientée dans cette représentation, et réciproquement.
Pour des variations conjointes,
\begin{equation}
 d\eta_{\mathrm{ret}}
 =\frac{\eta_{\mathrm{ret}}}{\alpha}\,d\alpha
 -250\alpha\,\operatorname{sech}^2(\chi_{\mathrm{rad}}/2)
 \,d\chi_{\mathrm{rad}}.
 \label{md:restop:diferencial-accion}
\end{equation}

Le caractère $H_5$, la correction régionale et l’anisotropie radiale orientée
maintiennent ainsi une relation de compatibilité et ne constituent pas trois
paramètres indépendants. Le caractère analytique et la géométrie de
réciprocité représentent la même section de retour.
Ce niveau structurel de l’équation de Planck demeure explicite
avant de composer l’action avec une fréquence.

La décennie d’action conserve sa généalogie déterminantale, indépendante
de cette reparamétrisation :
\begin{equation}
 \hbar_{\mathrm{ret}}
 =10^{-34}\mathcal U_S\eta_{\mathrm{ret}},\qquad
 \hbar_{\mathrm{pre}}=10^{-34}\mathcal U_SH_5.
 \label{md:restop:secciones-accion}
\end{equation}
On conserve en outre l’action par tour $h=2\pi\hbar$ et la lecture
thermique
\begin{equation}
 k_B\Theta_{\mathrm{clk}}\log3=\frac{h}{T_{108}}
 \label{md:restop:lectura-termica}
\end{equation}
dans son domaine. L’horloge physique d’une réalisation conserve son constructeur :
la coordonnée $\chi_{\mathrm{rad}}$ classe cette réciprocité et non la
totalité des histoires TPK. La classification par
$\chi_{\mathrm{rad}}$ et la relation différentielle réunissent les représentations
existantes de l’action ; elles maintiennent leur caractère de formalisation focale et ne
constituent pas une seconde dérivation indépendante de Planck.

\subsection{Circuits de phase, mémoire et conservativité}
\label{md:restop:conservatividad}

L’holonomie nonadique restaure la phase et accroît la mémoire, selon
la section~\ref{mono-ampl:vueltas}. Si un parcours fermé dans les phases
augmente le compteur entier, cet accroissement ne peut pas s’écrire comme
différence d’un potentiel dépendant uniquement de la phase.

La démonstration s’obtient par télescopage : la somme
\begin{equation}
 \sum_j\bigl[F(g_{j+1})-F(g_j)\bigr]=0
 \label{md:restop:potencial-fase}
\end{equation}
sur un circuit de phases contraste avec l’accroissement de mémoire, égal
à un par tour. Sur l’état élargi, le compteur $m$ enregistre
cet accroissement. Le parcours fermé dans la projection conservait
donc une ouverture dans l’état complet.

L’étude de la conservativité commence par déclarer l’espace dans
lequel le parcours se ferme et les coordonnées qui participent à la
dynamique. Non-commutativité et non-conservativité sont des propriétés
différentes. La section~\ref{md:restop:memoria} contrôle l’ordre des
opérations ; cette section distingue un circuit projeté d’un circuit
complet. Pour étudier une force concrète, la fonctionnelle d’action ou de
travail doit agir sur le même espace d’états et le même transport.

La croissance de la mémoire n’attribue pas à elle seule une production
thermodynamique. La section~\ref{md:cron:crecimiento} démontre que
l’information des blocs peut croître avec un taux entropique nul.
La restitution opératorielle ajoute une précision : conserver l’information
peut signifier conserver les transformations qui peuvent être composées,
en plus de maintenir les symboles stockés.

\subsection{Compatibilité dynamique entre observations}
\label{md:restop:sintesis}

L’unité d’étude est la compatibilité dynamique entre observations,
la mémoire faisant partie de l’opération. La valeur d’une constante
constitue une lecture ; sa structure contient des contraintes sur d’autres
lectures et sur les transformations admissibles. Le prolongement de
l’analyse consiste à obtenir des relations qui empêchent de modifier de manière
indépendante la valeur, l’orientation, la mémoire, la forme et l’action.

Les trois développements focaux parcourent ce même axe :
\begin{itemize}
 \item Un cycle complet de régimes TRIT conserve encore un
 déplacement nonadique de capacité : la régularité locale n’épuise pas
 le retour.
 \item La réentrée de mémoire modifie exactement la réponse à l’ordre
 de deux opérations : l’historique intervient causalement dans la composition.
 \item Le caractère $H_5$ et la géométrie radiale orientée déterminent
 mutuellement leur section de retour : l’action admet une représentation
 structurelle récupérable.
\end{itemize}
Ces conséquences prolongent la généalogie de $\pi,\varphi,e,\alpha$ et
explicitent ce que signifie la conserver lors du passage aux opérations et à
l’action. La composition ultérieure réunit ces niveaux au moyen des applications
communes déjà construites, en maintenant la distinction entre leurs indices et
leurs mémoires même s’ils partagent une dénomination.
